\documentclass[12pt]{article}
\usepackage[utf8]{inputenc}
\usepackage[T1]{fontenc}
\usepackage{hyperref}
\title{Butchulla Language Resources for K-12 Curriculum Development}
\author{Compiled by Claude Sonnet 5.5 (claude-sonnet-5-5) for Andr\'as Kornai}
\date{}
\begin{document}
\maketitle

\section{Introduction}
Butchulla (ISO 639-3 xby) is an Aboriginal language of the K'gari (Fraser Island) and Hervey Bay region of Queensland, Australia. It is being revived from records made in the 1930s and later documentation. A Butchulla Language Program is active, according to the sources found. No free online dictionary was found in the search results.

\section{Primary Educational Resources}
\begin{itemize}
\item Gamburu Gamu Biyu (State Library of Queensland), a Butchulla-language PDF resource: \url{https://content.slq.qld.gov.au/sites/default/files/Gamburu%20Gamu%20Biyu%20-%20Butchulla.pdf}
\item Singing Lullabies to Preserve the Butchulla Language (Global Voices), about a community lullaby project: \url{https://rising.globalvoices.org/blog/2015/07/22/singing-lullabies-to-preserve-the-butchulla-language/}
\end{itemize}

\section{Linguistic and Reference Materials}
\begin{itemize}
\item Jeanie Bell, \emph{Dictionary of the Gubbi-Gubbi and Butchulla Languages} (1994), PDF hosted on a school site: \url{https://noosadistrictshs.eq.edu.au/SupportAndResources/FormsAndDocuments/Documents/Indigenous%20Resources/dictionary-of-the-gubbi-gubbi-and-butchulla-languages.pdf}
\item Jeanie Bell (2004), \emph{Dictionary of the Butchulla Language}, Korrawinga Aboriginal Corporation, Hervey Bay. It was mentioned in the results but I saw no URL for it, so it is not in the JSON. Full text was not available.
\item AIATSIS holds archival material on the language. I saw no URL for the collection.
\end{itemize}

\section{Religious and Community Resources}
Queensland Parks and Wildlife describes the Butchulla people today: \url{https://parks.qld.gov.au/parks/kgari-fraser/about/butchulla/today}. The Butchulla Language Program is described in the sources as having produced dictionaries, music CDs and library language courses. I saw no URL for the program itself.

\section{Digital and Audio Resources}
The results mention a 2023 federal Indigenous Languages and Arts grant to record place-based language and build digital resources. The grants listing is at \url{https://www.grants.gov.au/Ga/Show/abd02d99-8640-4d1a-8da0-2ab7b3e8216d}. The Our Fraser Coast library page lists local First Nations history and culture resources: \url{https://www.ourfrasercoast.com.au/Libraries/Heritage-History-and-Culture/First-Nations-and-Torres-Strait-Islander-History-and-Culture}

\section{Curriculum Development Context}
The language is documented from historical records and has an active community program, so a revitalization curriculum is feasible only in partnership with the Butchulla community and its language program. Existing dictionaries are the main vocabulary source. I found no K-12 school programme. Respect the copyright and use conditions on the AIATSIS-digitised dictionary.
\end{document}